\documentclass{standalone}
\usepackage{tikz}
\usepackage{animate}
\usepackage{graphicx}
\usetikzlibrary{fit, positioning, arrows.meta, shapes, calc}

\begin{document}

% Macro for drawing 3D Convolution Blocks
% Usage: \convblock{x}{y}{width}{height}{depth}
\def\convblock#1#2#3#4#5{
    \pgfmathsetmacro{\w}{#3/2}
    \pgfmathsetmacro{\h}{#4/2}
    \pgfmathsetmacro{\dx}{#5*0.6}
    \pgfmathsetmacro{\dy}{#5*0.4}
    
    % Top face
    \filldraw[fill=orange!50!brown, draw=brown!80!black, thick, line join=round]
        (#1-\w, #2+\h) -- (#1+\w, #2+\h) -- (#1+\w+\dx, #2+\h+\dy) -- (#1-\w+\dx, #2+\h+\dy) -- cycle;
    % Right face
    \filldraw[fill=orange!70!brown, draw=brown!80!black, thick, line join=round]
        (#1+\w, #2-\h) -- (#1+\w, #2+\h) -- (#1+\w+\dx, #2+\h+\dy) -- (#1+\w+\dx, #2-\h+\dy) -- cycle;
    % Front face
    \filldraw[fill=orange!40!brown, draw=brown!80!black, thick, line join=round]
        (#1-\w, #2-\h) rectangle (#1+\w, #2+\h);
}

\begin{animateinline}[controls, loop]{1}
\multiframe{14}{iFrame=1+1}{

\pgfmathsetmacro{\MaxFrame}{14}

% Opacity Logic
% Step 1: Image
\pgfmathsetmacro{\opImage}{\iFrame==1 || \iFrame==\MaxFrame ? 1 : 0.15}
% Step 2: Offline LiDAR map
\pgfmathsetmacro{\opMap}{\iFrame==2 || \iFrame==\MaxFrame ? 1 : 0.15}
% Step 3: Projected depth
\pgfmathsetmacro{\opDepth}{\iFrame==3 || \iFrame==\MaxFrame ? 1 : 0.15}
% Step 4: Image Encoder
\pgfmathsetmacro{\opImgEnc}{\iFrame==4 || \iFrame==\MaxFrame ? 1 : 0.15}
% Step 5: LiDAR Encoder
\pgfmathsetmacro{\opLidarEnc}{\iFrame==5 || \iFrame==\MaxFrame ? 1 : 0.15}
% Step 6: Correlation
\pgfmathsetmacro{\opCorr}{\iFrame==6 || \iFrame==\MaxFrame ? 1 : 0.15}
% Step 7: 3D Cost volume
\pgfmathsetmacro{\opCostVol}{\iFrame==7 || \iFrame==\MaxFrame ? 1 : 0.15}
% Step 8: Position embeddings
\pgfmathsetmacro{\opPosEmb}{\iFrame==8 || \iFrame==\MaxFrame ? 1 : 0.15}
% Step 9: Element-wise add
\pgfmathsetmacro{\opPlus}{\iFrame==9 || \iFrame==\MaxFrame ? 1 : 0.15}
% Step 10: POET
\pgfmathsetmacro{\opPoet}{\iFrame==10 || \iFrame==\MaxFrame ? 1 : 0.15}
% Step 11: Relative pose
\pgfmathsetmacro{\opRelPose}{\iFrame==11 || \iFrame==\MaxFrame ? 1 : 0.15}
% Step 12: Multiply
\pgfmathsetmacro{\opMult}{\iFrame==12 || \iFrame==\MaxFrame ? 1 : 0.15}
% Step 13: Optimized vehicle pose
\pgfmathsetmacro{\opOptPose}{\iFrame==13 || \iFrame==\MaxFrame ? 1 : 0.15}

% Parent Container Opacity (Main Dashed Box)
% Must be visible from Step 4 (first internal element) through Step 13 (last internal element) and Finale
\pgfmathsetmacro{\opContainer}{(\iFrame>=4 && \iFrame<=13) || \iFrame==\MaxFrame ? 1 : 0.15}

\begin{tikzpicture}[
    >=Stealth,
    label_text/.style={
        fill=white,
        inner sep=2pt,
        font=\bfseries\itshape
    },
    image_node/.style={
        inner sep=0pt
    }
]

\path [use as bounding box] (-4, -7) rectangle (27, 10);

% -----------------------------------------------------
% 1. Input Image Nodes (Left Side)
% -----------------------------------------------------
\begin{scope}[opacity=\opImage, transparency group]
    \node[image_node] (image) at (0, 4) {\includegraphics[width=4cm, height=2cm, keepaspectratio]{crops/2305.04298v1-Figure1-1/image.png}};
    \node[below=0.1cm of image, font=\sffamily] {Image};
    % Arrow to Map (part of input flow)
    \draw[->, thick] (image.west) -- ++(-1.0, 0) |- (map.west);
    \node[right] at (-3.0, 2) {$\mathcal{P}_o$};
\end{scope}

\begin{scope}[opacity=\opMap, transparency group]
    \node[image_node] (map) at (0, 0) {\includegraphics[width=3cm, height=2.5cm, keepaspectratio]{crops/2305.04298v1-Figure1-1/map.png}};
    \node[right=0.2cm of map, align=left, font=\sffamily] {Offline\\LiDAR map};
    % Arrow to Depth
    \draw[->, thick] (map.south) -- (depth.north);
\end{scope}

\begin{scope}[opacity=\opDepth, transparency group]
    \node[image_node] (depth) at (0, -4) {\includegraphics[width=4cm, height=2cm, keepaspectratio]{crops/2305.04298v1-Figure1-1/depth.png}};
    \node[below=0.1cm of depth, font=\sffamily] {Projected depth};
\end{scope}

% -----------------------------------------------------
% 2. Image Encoder (Top Branch)
% -----------------------------------------------------
\begin{scope}[opacity=\opImgEnc, transparency group]
    \convblock{3.6}{4}{0.3}{3.0}{0.8}
    \convblock{4.2}{4}{0.3}{2.4}{0.8}
    \convblock{4.8}{4}{0.4}{1.8}{0.8}
    \convblock{5.4}{4}{0.4}{1.4}{0.8}
    \convblock{6.0}{4}{0.5}{1.0}{0.8}
    \convblock{6.6}{4}{0.5}{0.6}{0.8}
    \coordinate (enc1_tl) at (3.4, 5.8);
    \coordinate (enc1_br) at (7.3, 2.2);
    \node[fit=(enc1_tl) (enc1_br)] (image_encoder) {};
    \node[label_text] at (5.1, 4) {Image Encoder};
    % Connection from Image to Encoder
    \draw[->, thick] (image.east) -- (image_encoder.west |- image.east);
\end{scope}

% -----------------------------------------------------
% 3. LiDAR Encoder (Bottom Branch)
% -----------------------------------------------------
\begin{scope}[opacity=\opLidarEnc, transparency group]
    \convblock{3.6}{-4}{0.3}{3.0}{0.8}
    \convblock{4.2}{-4}{0.3}{2.4}{0.8}
    \convblock{4.8}{-4}{0.4}{1.8}{0.8}
    \convblock{5.4}{-4}{0.4}{1.4}{0.8}
    \convblock{6.0}{-4}{0.5}{1.0}{0.8}
    \convblock{6.6}{-4}{0.5}{0.6}{0.8}
    \coordinate (enc2_tl) at (3.4, -2.2);
    \coordinate (enc2_br) at (7.3, -5.8);
    \node[fit=(enc2_tl) (enc2_br)] (lidar_encoder) {};
    \node[label_text] at (5.1, -4) {LiDAR Encoder};
    % Connection from Depth to Encoder
    \draw[->, thick] (depth.east) -- (lidar_encoder.west |- depth.east);
\end{scope}

% -----------------------------------------------------
% 4. Correlation Block
% -----------------------------------------------------
\begin{scope}[opacity=\opCorr, transparency group]
    \node[draw=violet!50!black, fill=violet!70!blue, text=white, thick, minimum width=1.0cm, minimum height=4.5cm, inner sep=0pt] (correlation) at (8.5, 0) {\rotatebox{-90}{\textbf{\textit{\Large Correlation}}}};
    % Connections from Encoders to Correlation
    \draw[->, thick] (image_encoder.east) -- ([yshift=1.5cm]correlation.west);
    \draw[->, thick] (lidar_encoder.east) -- ([yshift=-1.5cm]correlation.west);
\end{scope}

% -----------------------------------------------------
% 5. 3D Cost Volume
% -----------------------------------------------------
\begin{scope}[opacity=\opCostVol, transparency group]
    \pgfmathsetmacro{\cvx}{11.0}
    \pgfmathsetmacro{\cvy}{0}
    \pgfmathsetmacro{\cvw}{1.3}
    \pgfmathsetmacro{\cvh}{1.0}
    \pgfmathsetmacro{\cvdx}{0.8}
    \pgfmathsetmacro{\cvdy}{0.6}

    % Back face
    \filldraw[fill=red!5!white, draw=red!20, thick, opacity=0.6] (\cvx-\cvw+\cvdx, \cvy-\cvh+\cvdy) rectangle (\cvx+\cvw+\cvdx, \cvy+\cvh+\cvdy);
    % Inner heat lines
    \draw[red!60, line width=1.5pt, opacity=0.7] (\cvx-\cvw+0.5+\cvdx, \cvy-\cvh+0.5+\cvdy) -- (\cvx+\cvw-0.5+\cvdx, \cvy-\cvh+0.5+\cvdy) -- (\cvx+\cvw-0.5+\cvdx, \cvy+\cvh-0.5+\cvdy);
    \draw[red!80, line width=2pt, opacity=0.8] (\cvx-\cvw+0.3, \cvy-\cvh+0.3) -- (\cvx+\cvw-0.3, \cvy-\cvh+0.3) -- (\cvx+\cvw-0.3, \cvy+\cvh-0.3);
    \fill[red!40, opacity=0.5] (\cvx, \cvy-\cvh+0.1) rectangle (\cvx+0.5, \cvy+\cvh-0.2);
    % Bottom face
    \filldraw[fill=red!10!white, draw=red!20, thick, opacity=0.6] (\cvx-\cvw, \cvy-\cvh) -- (\cvx+\cvw, \cvy-\cvh) -- (\cvx+\cvw+\cvdx, \cvy-\cvh+\cvdy) -- (\cvx-\cvw+\cvdx, \cvy-\cvh+\cvdy) -- cycle;
    % Top face
    \filldraw[fill=red!10!white, draw=red!20, thick, opacity=0.6] (\cvx-\cvw, \cvy+\cvh) -- (\cvx+\cvw, \cvy+\cvh) -- (\cvx+\cvw+\cvdx, \cvy+\cvh+\cvdy) -- (\cvx-\cvw+\cvdx, \cvy+\cvh+\cvdy) -- cycle;
    % Right face
    \filldraw[fill=red!10!white, draw=red!20, thick, opacity=0.6] (\cvx+\cvw, \cvy-\cvh) -- (\cvx+\cvw, \cvy+\cvh) -- (\cvx+\cvw+\cvdx, \cvy+\cvh+\cvdy) -- (\cvx+\cvw+\cvdx, \cvy-\cvh+\cvdy) -- cycle;
    % Front face
    \filldraw[fill=red!15!white, draw=red!30, thick, opacity=0.6] (\cvx-\cvw, \cvy-\cvh) rectangle (\cvx+\cvw, \cvy+\cvh);

    \coordinate (cv_tl) at (\cvx-\cvw-0.2, \cvy+\cvh+\cvdy+0.2);
    \coordinate (cv_br) at (\cvx+\cvw+\cvdx+0.2, \cvy-\cvh-0.2);
    \node[fit=(cv_tl) (cv_br)] (cost_volume) {};
    \node[below] at (\cvx+0.4, \cvy-\cvh-0.1) {3D Cost volume};
    
    % Connection Correlation to Cost Volume
    \draw[->, thick] (correlation.east) -- (cost_volume.west |- correlation.east);
\end{scope}

% -----------------------------------------------------
% 6. Position Embeddings
% -----------------------------------------------------
\begin{scope}[opacity=\opPosEmb, transparency group]
    \pgfmathsetmacro{\pex}{11.0}
    \pgfmathsetmacro{\pey}{4.0}
    \foreach \i in {4,3,2,1,0} {
        \filldraw[fill=gray!15, draw=gray!60, thick] (\pex - 0.8 + \i*0.15, \pey - 1.0 + \i*0.15) rectangle (\pex + 0.8 + \i*0.15, \pey + 1.0 + \i*0.15);
    }
    \coordinate (pe_tl) at (\pex - 1.0, \pey + 1.8);
    \coordinate (pe_br) at (\pex + 1.6, \pey - 1.2);
    \node[fit=(pe_tl) (pe_br)] (pos_embeddings) {};
    \node[above] at (\pex+0.4, \pey+1.6) {Position embeddings};
\end{scope}

% -----------------------------------------------------
% 7. Element-wise Add (Plus) Node
% -----------------------------------------------------
\begin{scope}[opacity=\opPlus, transparency group]
    \node[draw=black, circle, thick, inner sep=0pt, minimum size=0.6cm, fill=white] (plus) at (14.0, 4.0) {\Large $+$};
    % Connections to Plus
    \draw[->, thick] (pos_embeddings.east |- plus.west) -- (plus.west);
    \draw[->, thick] (cost_volume.east) -- (plus.south);
\end{scope}

% -----------------------------------------------------
% 8. POET Block
% -----------------------------------------------------
\begin{scope}[opacity=\opPoet, transparency group]
    \node[draw=black, thick, fill=cyan!10, minimum width=3.8cm, minimum height=2.4cm, align=center] (poet) at (17.5, 4.0) {
        \textbf{\textit{\textcolor{red}{\Large POET}}} \\[0.2cm] 
        \textbf{\textcolor{red}{PO}}se \textbf{\textcolor{red}{E}}stimator \\ 
        \textbf{\textcolor{red}{T}}ransformer
    };
    % POET self-loop
    \draw[->, thick] ([xshift=1cm]poet.north) -- ++(0, 0.8) -| ([xshift=-0.5cm]poet.north);
    % Connection Plus to POET
    \draw[->, thick] (plus.east) -- (poet.west);
\end{scope}

% -----------------------------------------------------
% 9. Relative Pose (Frustums)
% -----------------------------------------------------
\begin{scope}[opacity=\opRelPose, transparency group]
    % Left Frustum
    \draw[thick, line join=round] (21.4, 4.4) -- (21.8, 4.7) -- (21.8, 3.7) -- (21.4, 3.4) -- cycle;
    \draw[thick] (22.2, 3.8) -- (21.4, 4.4);
    \draw[thick] (22.2, 3.8) -- (21.8, 4.7);
    \draw[thick] (22.2, 3.8) -- (21.8, 3.7);
    \draw[thick] (22.2, 3.8) -- (21.4, 3.4);
    \draw[thin] (21.6, 4.55) -- (21.6, 3.55); % camera plane detail

    % Right Frustum
    \draw[thick, line join=round] (24.0, 4.4) -- (23.6, 4.7) -- (23.6, 3.7) -- (24.0, 3.4) -- cycle;
    \draw[thick] (23.2, 3.8) -- (24.0, 4.4);
    \draw[thick] (23.2, 3.8) -- (23.6, 4.7);
    \draw[thick] (23.2, 3.8) -- (23.6, 3.7);
    \draw[thick] (23.2, 3.8) -- (24.0, 3.4);
    \draw[thin] (23.8, 4.55) -- (23.8, 3.55); % camera plane detail

    % Double arrow
    \draw[<->, thick] (22.3, 3.9) to[bend right=25] node[below] {$\mathcal{P}_\Delta$} (23.1, 3.9);

    % Fit group for Relative Pose
    \coordinate (rp_tl) at (21.2, 5.0);
    \coordinate (rp_br) at (24.2, 2.8);
    \node[fit=(rp_tl) (rp_br)] (rel_pose) {};
    \node[below] at (22.7, 3.2) {Relative pose};
    
    % Connection POET to Relative Pose
    \draw[->, thick] (poet.east) -- (rel_pose.west |- poet.east);
\end{scope}

% -----------------------------------------------------
% 10. Multiply Node & Outputs
% -----------------------------------------------------
\begin{scope}[opacity=\opMult, transparency group]
    \node[draw=black, circle, thick, inner sep=0pt, minimum size=0.6cm, fill=white] (multiply) at (22.7, 0.5) {\Large $\times$};
    % Connection Relative Pose to Multiply
    \draw[->, thick] (rel_pose.south) -- (multiply.north);
    % P_o input to Multiply
    \draw[<-, thick] (multiply.east) -- ++(1.5, 0) node[right] {$\mathcal{P}_o$};
\end{scope}

\begin{scope}[opacity=\opOptPose, transparency group]
    % Optimized vehicle pose
    \node[image_node] (opt_pose) at (22.7, -3.0) {\includegraphics[width=3.5cm, height=2cm, keepaspectratio]{crops/2305.04298v1-Figure1-1/final_pose.png}};
    \node[below=0.1cm of opt_pose, align=center, font=\sffamily] {Optimized\\vehicle pose};
    % Connection Multiply to Optimized Pose
    \draw[->, thick] (multiply.south) -- (opt_pose.north);
\end{scope}

% -----------------------------------------------------
% 11. Main Dashed Bounding Box & Legend
% -----------------------------------------------------
\begin{scope}[opacity=\opContainer, transparency group]
    \draw[thick, dashed] (2.5, 7.0) rectangle (20.5, -6.0);
    \node[fill=white, font=\Large\bfseries, inner sep=5pt] at (11.5, 7.0) {Image-to-LiDAR map localization network};
\end{scope}

% Legend Box (Always visible or part of container? Usually static. Let's make it part of container or always visible. 
% Given strict constraints, let's attach it to the container opacity or make it always visible. 
% The prompt says "Parent Container Persistence". The legend is outside the main dashed box. 
% I will make the legend always visible (opacity 1) as it explains the symbols, or attach to container.
% To be safe and clean, I'll attach it to the container opacity so it appears when the network appears.
\begin{scope}[opacity=\opContainer, transparency group]
    \draw[dashed, thick] (14.5, -3.5) rectangle (19.5, -5.5);
    \convblock{15.2}{-4.0}{0.2}{0.6}{0.4}
    \node[right, font=\normalsize] at (15.8, -4.0) {Convolution block};

    \node[draw=black, circle, thick, inner sep=0pt, minimum size=0.5cm] (leg_plus) at (15.2, -4.9) {$+$};
    \node[right, font=\normalsize] at (15.8, -4.9) {Element-wise add};
\end{scope}

\end{tikzpicture}
}
\end{animateinline}

\end{document}